% !TEX program = xelatex
% ============================================================
%  ملاحظات المحاضرة — Lecture Notes Template
%  قالب ملاحظات محاضرة مع ترويسة (المقرر، رقم المحاضرة، التاريخ، المعلم)
%  ومخطط مواضيع، مفاهيم رئيسية في صناديق ملوّنة، تعريفات ونظريات،
%  أمثلة، خلاصة، ومراجع.
%  Compile with: xelatex main.tex (run twice)
% ============================================================
%  Features:
%    - A4, 11pt, article class
%    - Header block: course, lecture #, date, instructor
%    - Topics outline, key concepts (tcolorbox callouts)
%    - Definitions / Theorems (amsthm-style custom environments)
%    - Examples, Summary, References
%    - Section dividers, fancyhdr header/footer
% ============================================================

\documentclass[11pt,a4paper]{article}

% ---- Geometry ----
\usepackage[a4paper,margin=2.5cm,top=2.5cm,bottom=2.5cm,headheight=16pt]{geometry}

% ---- Core packages ----
\usepackage{graphicx}
\usepackage{array}
\usepackage{booktabs}
\usepackage{tabularx}
\usepackage{multirow}
\usepackage{enumitem}
\usepackage{float}
\usepackage{titlesec}
\usepackage{fancyhdr}
\usepackage{setspace}
\usepackage{xcolor}
\usepackage{amsmath}
\usepackage{tcolorbox}

% ---- RTL / Arabic language support (order matters) ----
\usepackage{bidi}
\usepackage[no-math]{fontspec}
\usepackage[quiet]{polyglossia}

% ---- Fonts ----
\setmainfont[Scale=1]{Amiri-Regular.ttf}
\newfontfamily\arabicfont[Script=Arabic,Scale=0.95]{Amiri-Regular.ttf}
\newfontfamily\englishfont[Scale=0.95]{TeX Gyre Termes}

% ---- Languages ----
\setdefaultlanguage[calendar=gregorian,hijricorrection=1]{arabic}
\setotherlanguage[variant=british]{english}

% ---- Hyperref (load last, after polyglossia/bidi) ----
\usepackage[breaklinks,hidelinks]{hyperref}

% ---- Colors ----
\definecolor{accent}{HTML}{1F3A5F}
\definecolor{conceptbg}{HTML}{EAF1F8}
\definecolor{conceptframe}{HTML}{1F3A5F}
\definecolor{defbg}{HTML}{F0F7EE}
\definecolor{defframe}{HTML}{2E7D32}
\definecolor{thmbg}{HTML}{FBF3E4}
\definecolor{thmframe}{HTML}{B26A00}
\definecolor{exbg}{HTML}{F4F0FA}
\definecolor{exframe}{HTML}{6A1B9A}

% ---- Line spacing ----
\setstretch{1.4}

% ---- Captions in Arabic ----
\addto\captionsarabic{%
  \renewcommand{\figurename}{الشكل}%
  \renewcommand{\tablename}{الجدول}%
  \renewcommand{\refname}{المراجع}%
}

% ---- Section formatting ----
\titleformat{\section}{\Large\bfseries\color{accent}}{\thesection}{0.7em}{}
\titleformat{\subsection}{\large\bfseries}{\thesubsection}{0.6em}{}
\titlespacing*{\section}{0pt}{1.2ex plus 0.3ex}{0.8ex plus 0.2ex}

% ---- Section divider rule ----
\newcommand{\sectiondivider}{%
  \par\noindent{\color{accent}\rule{\linewidth}{0.8pt}}\par\vspace{0.4em}%
}

% ---- Headers / footers ----
\pagestyle{fancy}
\fancyhf{}
\renewcommand{\headrulewidth}{0.4pt}
\renewcommand{\footrulewidth}{0pt}
% TODO: عدّل اسم المقرر ورقم المحاضرة في الترويسة.
\fancyhead[R]{\small ملاحظات المحاضرة (Lecture Notes)}
\fancyhead[L]{\small اسم المؤسسة (Institution Name)}
\fancyfoot[C]{\small \thepage}

% ============================================================
% tcolorbox-based callout environments
% ============================================================

% Key concept box
\newtcolorbox{conceptbox}[1]{%
  colback=conceptbg, colframe=conceptframe,
  boxrule=0.6pt, arc=2pt, left=8pt, right=8pt, top=4pt, bottom=4pt,
  title={#1}, fonttitle=\bfseries\color{white},
  colbacktitle=conceptframe
}

% Definition box
\newtcolorbox{definitionbox}[1]{%
  colback=defbg, colframe=defframe,
  boxrule=0.6pt, arc=2pt, left=8pt, right=8pt, top=4pt, bottom=4pt,
  title={تعريف: #1}, fonttitle=\bfseries\color{white},
  colbacktitle=defframe
}

% Theorem box
\newtcolorbox{theoremtbox}[1]{%
  colback=thmbg, colframe=thmframe,
  boxrule=0.6pt, arc=2pt, left=8pt, right=8pt, top=4pt, bottom=4pt,
  title={نظرية: #1}, fonttitle=\bfseries\color{white},
  colbacktitle=thmframe
}

% Example box
\newtcolorbox{examplebox}[1]{%
  colback=exbg, colframe=exframe,
  boxrule=0.6pt, arc=2pt, left=8pt, right=8pt, top=4pt, bottom=4pt,
  title={مثال: #1}, fonttitle=\bfseries\color{white},
  colbacktitle=exframe
}

\begin{document}

% ============================================================
% TITLE / HEADER BLOCK
% ============================================================
\begin{center}
{\LARGE\bfseries ملاحظات المحاضرة (Lecture Notes)}\\[0.4em]
{\large\bfseries اسم المؤسسة (Institution Name)}
\end{center}

\vspace{0.3em}
{\color{accent}\rule{\linewidth}{1.2pt}}
\vspace{0.5em}

% ---- Info table ----
\noindent
\begin{tabularx}{\linewidth}{@{}p{0.28\linewidth}X@{}}
\textbf{المقرر (Course):} & % TODO: ضع اسم المقرر هنا
اسم المقرر (Course Name) \\
\textbf{رقم المحاضرة (Lecture \#):} & % TODO: ضع رقم المحاضرة هنا
\hrulefill \\
\textbf{التاريخ (Date):} & % TODO: ضع التاريخ هنا
\hrulefill \\
\textbf{اسم المعلم (Instructor):} & % TODO: ضع اسم المعلم هنا
\hrulefill \\
\end{tabularx}

\vspace{0.4em}
{\color{accent}\rule{\linewidth}{0.8pt}}
\vspace{0.6em}

% ============================================================
% TOPICS OUTLINE
% ============================================================
\section{مخطط المواضيع (Topics Outline)}
\sectiondivider
% TODO: عدّل قائمة المواضيع حسب محتوى المحاضرة.
\begin{itemize}[leftmargin=2em, itemsep=0.2em]
  \item الموضوع الأول (Topic One)
  \item الموضوع الثاني (Topic Two)
  \item الموضوع الثالث (Topic Three)
  \item الموضوع الرابع (Topic Four)
\end{itemize}

% ============================================================
% KEY CONCEPTS
% ============================================================
\section{المفاهيم الرئيسية (Key Concepts)}
\sectiondivider
% TODO: استبدل المفاهيم أدناه بمفاهيم المحاضرة الفعلية.

\begin{conceptbox}{المفهوم الأول (Concept One)}
وصف موجز للمفهوم الأول وأهميته في سياق المحاضرة. اكتب هنا النقاط الأساسية
التي يجب على الطالب تذكّرها.
\end{conceptbox}

\vspace{0.5em}

\begin{conceptbox}{المفهوم الثاني (Concept Two)}
وصف موجز للمفهوم الثاني مع توضيح العلاقة بينه وبين المفاهيم الأخرى.
\end{conceptbox}

% ============================================================
% DEFINITIONS AND THEOREMS
% ============================================================
\section{التعريفات والنظريات (Definitions \& Theorems)}
\sectiondivider
% TODO: أضف التعريفات والنظريات الخاصة بالمحاضرة.

\begin{definitionbox}{التعريف الأول (Definition One)}
% TODO: ضع نص التعريف هنا.
نص التعريف الأول يُكتب هنا بشكل دقيق ومختصر.
\end{definitionbox}

\vspace{0.5em}

\begin{theoremtbox}{النظرية الأولى (Theorem One)}
% TODO: ضع نص النظرية هنا.
نص النظرية الأولى مع شروطها ونتيجتها الرئيسية.
\end{theoremtbox}

% ============================================================
% EXAMPLES
% ============================================================
\section{أمثلة (Examples)}
\sectiondivider
% TODO: استبدل الأمثلة أدناه بأمثلة المحاضرة الفعلية.

\begin{examplebox}{مثال 1 (Example 1)}
% TODO: ضع نص المثال والحل هنا.
نص المثال: احسب مشتقة الدالة $f(x) = x^{2} + 3x - 5$.

\textbf{الحل:}
\[
f'(x) = 2x + 3
\]
\end{examplebox}

\vspace{0.5em}

\begin{examplebox}{مثال 2 (Example 2)}
% TODO: ضع نص المثال والحل هنا.
نص المثال الثاني يُكتب هنا مع خطوات الحل بالتفصيل.
\end{examplebox}

% ============================================================
% SUMMARY
% ============================================================
\section{الخلاصة (Summary)}
\sectiondivider
% TODO: اكتب خلاصة المحاضرة هنا.
تلخّص هذه المحاضرة المفاهيم الأساسية المتعلّقة بالموضوع، وتطرّقت إلى
عدد من التعريفات والنظريات مع أمثلة تطبيقية. يجب على الطالب مراجعة
النقاط الرئيسية وحل التمارين المرفقة.

% ============================================================
% REFERENCES
% ============================================================
\section{المراجع (References)}
\sectiondivider
% TODO: أضف المراجع الفعلية هنا.
\begin{itemize}[leftmargin=2em, itemsep=0.2em]
  \item المرجع الأول (Reference One) — اسم المؤلف، \emph{عنوان الكتاب}، دار النشر، السنة.
  \item المرجع الثاني (Reference Two) — اسم المؤلف، \emph{عنوان الكتاب}، دار النشر، السنة.
\end{itemize}

\end{document}
